\section{Baselines}
\label{sec:baselines}

A benchmark that only ever scored one system would be a validation
appendix. Astronomy~v1 ships four reference baselines, chosen to
bracket the interesting comparisons; each consumes only \corpusA{}
records, freezes its output before any future-corpus access, and is
evaluated through the identical pipeline.

\begin{description}
  \item[B1: Random claims.] Sample random pre-cutoff papers and template
    their headline result into a robustness question. The floor: any
    system must beat chance-directed attention.
  \item[B2: Direct LLM.] Give an LLM a sample of pre-cutoff abstracts
    and ask for the most valuable open questions. The ``why not just ask
    GPT?'' comparison. Because a modern LLM's weights postdate the
    cutoff, this baseline is also a contamination probe: performance
    that vanishes for post-training-cutoff instances indicates
    memorized hindsight rather than generation ability
    (Section~\ref{sec:threats}).
  \item[B3: Review future work.] Extract explicitly posed open questions
    from pre-cutoff review papers. The strongest natural reference: a
    discovery system is interesting only if it adds value over questions
    the community had already written down. Note this baseline should
    score highly on coverage \emph{by construction}---these questions
    are known community priorities---so the discriminating metrics are
    premise refutation and lead time, where a copied question can never
    lead the field.
  \item[B4: Citation leaders.] Template follow-up questions from the
    most-cited pre-cutoff papers. Tests whether chasing prominence
    matches structured evidence analysis.
\end{description}

\begin{table}[t]
  \centering
  \small
  \caption{Astronomy v1 leaderboard. The pilot release scores one
  submission; baseline rows will be populated from the released harness
  (\texttt{baselines/}), and external systems can submit via the
  interface in Section~\ref{sec:protocol}. Lead time is reported as null
  until community first-posed dates are annotated
  (Section~\ref{sec:metrics}).}
  \label{tab:leaderboard}
  \begin{tabular}{@{}lccccc@{}}
    \toprule
    System & $n$ & Coverage & Answered & Premise refuted & Lead time \\
    \midrule
    \texttt{evidence\_graph\_v1} \citep{paper1} & 10 & 100\% & 20\% & 10\% & --- \\
    B1 random claims        & --- & --- & --- & --- & --- \\
    B2 direct LLM           & --- & --- & --- & --- & --- \\
    B3 review future work   & --- & --- & --- & --- & --- \\
    B4 citation leaders     & --- & --- & --- & --- & --- \\
    \bottomrule
  \end{tabular}
\end{table}

The pilot release implements the four baselines as runnable submission
generators but reports scores only for \texttt{evidence\_graph\_v1}
(Table~\ref{tab:leaderboard}): baseline outcome labels require the same
retrieval--judging--adjudication budget as any submission, and we
prefer an honest partially populated leaderboard over rushed labels.
Populating the baseline rows is the release's first roadmap item, and
the leaderboard file (\texttt{results/leaderboard.json}) is designed for
external submissions to append to.
